\chapter{Ход работы}
\label{ch:chap2}

\section{Получение метрик блоков СМО}

В MATLAB Simulink есть удобный инструментарий для снятия метрик с блоков СМО. У каждого блока свой набор метрик. Метрики для каждого 
блока изображены ниже

\begin{figure}[H]
    \centering
    \includegraphics[width=1.0\textwidth]{stats.png}
    \caption{Метрики для блоков СМО}
\end{figure}

После проставления галочек возле тех метрик, которые нам нужны, возле блоков появятся дополнительные выходы, на которые будут подаваться
подсчитанные метрики. Метрики можно выводить на блок Scope, чтобы визуализировать метрики на графики.

\begin{figure}[H]
    \centering
    \includegraphics[width=1.0\textwidth]{out_stats.png}
    \caption{Вывод метрик}
\end{figure}

\section{Анализ метрик генератора заявок}

\begin{figure}[H]
    \centering
    \includegraphics[width=1.0\textwidth]{gen_metrics.png}
    \caption{Метрики генератора заявок}
\end{figure}

Метрика d - кол-во заявок, которое покинуло блок, w - среднее время между генерациями заявок, pe - кол-во заявок, которое сгенерировано,
но находятся в системе (допустим, очередь уже полная и не может взять заявку, тогда эта заявка остается в блоке и ожидает освобождения 
очереди).

\section{Анализ метрик очереди заявок}

\begin{figure}[H]
    \centering
    \includegraphics[width=1.0\textwidth]{q_metrics.png}
    \caption{Метрики очереди}
\end{figure}

Здесь dq - кол-во вышедших заявок, n - текущее кол-во заявок в блоке, l - средняя длина очереди, wq - среднее время пребывания заявки в очереди.

\begin{figure}[H]
    \centering
    \includegraphics[width=1.0\textwidth]{q_ex.png}
    \caption{Метрики очереди}
\end{figure}

Эта метрика показывает кол-во заявок, исключенных из очереди принудительно, т.е сколько заявок было потеряно.

\section{Анализ метрик сервера}

\begin{figure}[H]
    \centering
    \includegraphics[width=1.0\textwidth]{s_metrics1.png}
    \caption{Метрики сервера}
\end{figure}

Здесь util - степень использования сервера, pes - кол-во заявок, готовых к обслуживанию, но пока не поступивших на обслуживание,
np - кол-во ожидающих заявок во внутренней очереди сервера, ws - среднее время обслуживания.

\begin{figure}[H]
    \centering
    \includegraphics[width=1.0\textwidth]{s_metrics2.png}
    \caption{Метрики сервера}
\end{figure}

ds - кол-во вышедших из свервера заявок.


\endinput
